% The spoken body of the homily for the First Sunday of Advent, Year B.
% Continuous preaching, ready to read aloud: no heading, no direction to the
% preacher, no citation inside the speech. The loci, the relation to the
% reviewed interpretations, the audience, the word count and the pace estimate
% stand in the terminal note, which no one reads out. The blank lines are the
% paragraph breaks a speaker recovers his place by, and the vertical spaces
% divide the speech into its movements: the command to all; the people who
% confess their sleep and call God Father; the petition for the will to
% watch, and the work given with it; Paul's church and God's fidelity; the
% one thing to do; and the watch the Preface names, ending on Paul's promise.
%
% Scripture is spoken in the Douay--Rheims, the leaf's study translation. The
% speech names that Bible once, before the first quoted verse whose wording
% the assembly will not have heard; the Gospel's last verse, quoted earlier,
% reads alike in it and in the Lectionary. No oration, Preface or antiphon is
% quoted in any Missal wording: the Collect, the Preface and the Prayer after
% Communion are described in part, and the Communion verse is named as the
% one the Missal gives, so nothing said depends on which chant is sung.
% Nothing assumes that the psalm or the Alleluia verse was sung in its
% Lectionary form, or that the solemn blessing is used.

On the Mount of Olives, across the valley from the Temple, four of the
disciples, Peter, James, John and Andrew, have asked Jesus privately when the
Temple's stones will fall. He has answered with a long speech, the last he
gives in Saint Mark's Gospel before the account of his Passion begins. Near
its end he tells them that no one knows the day or the hour except the
Father. Then comes the short parable we have just heard. A man goes abroad.
He leaves his house to his servants, gives each of them his own work, and
orders the doorkeeper to stay awake. He may come back in the evening, at
midnight, at cockcrow or in the morning. Those are the four watches into
which the night was divided, so the doorkeeper cannot count on any hour of
it. And the last sentence turns from the four men to everyone else: ``What I
say to you, I say to all: Watch.''

The four men who heard that command died long before the end of the world,
and so, most likely, will we. A bishop once wrote to Saint Augustine with
questions about the end of the world, and Augustine took up this verse in
his reply. It was said, he answers, to the four on the mountain, to those who
came after them, to us, and to all who will come after us until the Lord's
last coming. The day will not find all of them alive. The command still binds
every one of them, because ``in whatever state each one's own last day finds
him, in that state the last day of the world will take him.''

\bigskip

We have heard ``Watch'' every Advent of our lives, and most years we find
that we have slept. The first reading is the prayer of a people who know
they have failed God and say so. It is taken from a long prayer near the end
of the Book of Isaiah, and that prayer is spoken among ruins. The people say
that their enemies have trodden down God's sanctuary, and near its end they
speak of the house where their fathers praised him, now burnt with fire.
They call the Lord their Father and their Redeemer. They ask him why he has
let them stray from his ways. In the old Douay Bible, translated from the
Latin, the question is sharper: ``Why hast thou made us to err, O Lord, from
thy ways: why hast thou hardened our heart, that we should not fear thee?''
Saint Jerome did not take that to mean that God makes anyone sin. What looks
like the cause, he says, is God's patience: he waits for the sinner to be
saved and does not correct him at once, and the sinner, left uncorrected,
goes further.

Then the people cry to God to tear open the heavens and come down, and they
confess. God is angry, and they have sinned. Their justice is like a soiled
rag. They have fallen like leaves, and their sins carry them off like the
wind. And then comes a line that describes the sleep the Gospel warns
against: ``There is none that calleth upon thy name: that riseth up, and
taketh hold of thee.''

The people who say this go straight on to do what they have just said that
nobody does. They call on God's name, and the name they use is Father: ``And
now, O Lord, thou art our father, and we are clay: and thou art our maker,
and we all are the works of thy hands.''

Jerome, commenting on that verse, gives their reasoning. If we look at our
merits, he says, we can only despair; if we look at God's mercy, we dare to
pray, because he is our Father. Clay has no right to answer back to the
potter; but these people know that they are sons, and so they dare to go on,
in the verse after today's reading, and ask him not to be angry with them
beyond measure.

\bigskip

The psalm appointed after the reading is also the prayer of a defeated
people. They ask the Shepherd of Israel to stir up his might and come. The
refrain, which the psalm repeats three times, asks God to turn them and show
them his face, so that they will be saved. Saint Augustine, preaching on that
refrain, explained it in one sentence: we have turned away from God, and
unless he turns us, we will not be turned.

The Collect at the beginning of this Mass prayed in the same way. It asked
God to give his faithful the will to go out with just works to meet Christ as
he comes. Christ
commands us to stay awake, and the Church asks God for the will to do it.

Asking for the will does not mean waiting for it with folded hands. In the
parable the master gives his servants authority, each his own work, before he
leaves. Saint Bede, writing on this Gospel, says that this authority is the
grace of the Holy Spirit, given so that the faithful may serve in good works.
Then he comes to the last sentence, the one spoken to all, and says that
every one of us is bidden to guard the door of his own heart, so that the old
enemy cannot break in by suggesting evil, because each will give account to
God for himself.

\bigskip

Saint Paul wrote the second reading to a church he had founded at Corinth,
and the letter whose opening we heard goes on to rebuke that church: for its
factions, for a case of incest, for lawsuits against one another, for
disorder at the Lord's Supper itself. Yet he begins by thanking God for them.
They lack no grace, he says, as they wait for Christ to be revealed, and
Christ will strengthen them to the end, so that on his day no charge can be
brought against them. Then he gives his reason: ``God is faithful.''

Saint John Chrysostom asked how Paul could praise a church like that. One of
his answers is that the praise was never theirs: it belonged to the grace of
God. ``God is faithful,'' he explains, means that God is true; ``now if
true, what things He hath promised He will also perform.''

Paul says that God is faithful at the very beginning of the letter,
Chrysostom adds, so that after the rebukes to come the Corinthians would not
fall into despair. Where the calling fails, he says, the failure does not
come from the one who calls, but from those who start away from him.

\bigskip

Isaiah's people did not wait until they had made themselves just before they
spoke to God. They named their sleep, and in the same breath they called him
Father. The Missal gives for Communion today a verse of a psalm: ``the Lord
will give goodness: and our earth shall yield her fruit.'' Augustine,
preaching on that verse, spoke of the sinner who confesses, and said to him:
``Your sins, then, you can have of your own; good fruit you will not have
unless the one to whom you confess gives it.''

So before Christmas, do what Isaiah's people did. In the sacrament of
confession, name to God the place where you have been asleep, and call on
him as your Father.

\bigskip

In the Preface today the Church gives thanks for two comings of Christ: the
first in the lowliness of the flesh he took, which opened the way of
salvation, and the second in the glory of his majesty. And it uses the
Gospel's word of those who give thanks: keeping watch, it says, they dare to
wait for what was promised. The prayer after Communion says that God uses
these mysteries to teach us, even now, what to love and what to hold to. At
this altar he goes on forming in us the will to watch that we asked for at
the beginning of this Mass.

You may dare to wait on the same ground on which Isaiah's people dared to
pray and Paul's Corinthians could hope to stand without reproach. God is
faithful. He has called you into the fellowship of his Son, Jesus Christ,
who will strengthen you to the end, so that on his day no charge will stand
against you.
